\section{A Laurent model and a filtration by multipartitions}
\label{sec:laurent}

To complement the surjection of Proposition~\ref{prop:kernel-inclusion},
we seek an upper bound for the Hilbert series of $B_0$.  One loop at each
vertex cancels the minuscule tangent denominators, giving a Laurent model
for this calculation.

Throughout this Section the flavor torus is trivial, and the quiver has at
least one loop at every vertex of positive gauge dimension.  Delete the
vertices of gauge dimension zero, and choose one loop at every remaining
vertex.  Let $Q'$ be the quiver obtained by deleting these chosen loops.
Write
\[
 n_{ij}=\#Q'_1(i,j),\qquad c_{ij}=n_{ij}+n_{ji},\qquad
 C(\mathbf v)=\frac12\sum_{i,j}c_{ij}v_iv_j.
\]
In particular, $c_{ii}=2n_{ii}$.  We use the set
\[
 \Omega=\{(i,a):i\in Q_0,\ 1\leq a\leq v_i\}
\]
of gauge slots, the polynomial ring $P=\C[x_{ia}:(i,a)\in\Omega]$, and
$\mathcal W=\prod_iS_{v_i}$.  Let $\mathcal B'$ be the multiset consisting
of the weights $x_{jb}-x_{ia}$ of all arrows of $Q'$ and $w_i$ copies of
$-x_{ia}$ at each gauge slot.  The zero weights arising from diagonal
entries of loops can be omitted: their coweight pairing is always zero.

For $\lambda\in\Z^\Omega$, put
\[
 E_\beta(\lambda)=\bigl(-\langle\beta,\lambda\rangle\bigr)_+,
 \qquad
 r_\lambda=e^\lambda\prod_{\beta\in\mathcal B'}
                      \beta^{E_\beta(\lambda)},
\]
where $u_+=\max(u,0)$ and $e^\lambda e^\gamma=e^{\lambda+\gamma}$.
These expressions belong to $\C(x_{ia})[\Z^\Omega]$.  For a polynomial
$p\in P^{\mathcal W_\lambda}$, with
$\mathcal W_\lambda=\operatorname{Stab}_{\mathcal W}(\lambda)$, denote
the distinct-translate orbit sum by
\[
 \mathsf S_\lambda(p)=
 \sum_{w\in\mathcal W/\mathcal W_\lambda}w(p r_\lambda).
\]
The choice of representative of the coweight orbit does not affect the
space of such sums.

\begin{proposition}\label{prop:laurent-model}
Under the abelianization embedding of
Proposition~\ref{prop:minuscule}, the no-flavor Coulomb algebra is
\begin{equation}\label{eq:laurent-model}
 A_0=\left(\bigoplus_{\lambda\in\Z^\Omega}P r_\lambda\right)^{\mathcal W}.
\end{equation}
The identification preserves the coefficient ring $R_0=P^{\mathcal W}$
and the topological grading.
\end{proposition}

\begin{proof}
First, the direct sum on the right is an algebra.  Indeed,
\begin{equation}\label{eq:laurent-product}
 r_\alpha r_\gamma
 =r_{\alpha+\gamma}
  \prod_{\beta\in\mathcal B'}
   \beta^{E_\beta(\alpha)+E_\beta(\gamma)-E_\beta(\alpha+\gamma)},
\end{equation}
and every exponent is nonnegative.  The multiset $\mathcal B'$ is
$\mathcal W$-stable, so taking invariants gives a subalgebra.

For a positive or negative subset coweight, the weights of the deleted
loop at a vertex cancel the tangent Euler denominator in
Equation~\eqref{eq:minuscule}, up to a nonzero constant sign.  This is the
Grassmannian tangent-weight calculation: for a positive subset $I$,
the loop contributes
$\prod_{a\in I,\ b\notin I}(x_{ib}-x_{ia})$, while the tangent weights
are these roots, up to their common sign convention.  For a negative
subset the roots are reversed.  This is the local Euler-factor calculation
associated with minuscule flag localization; compare the pushforward
identity in \cite[Lemma~2.10]{Weekes19}.
Proposition~\ref{prop:minuscule} therefore
places every generator of $A_0$ in the right-hand side of
Equation~\eqref{eq:laurent-model}.  Conversely, it places every
$\mathsf S_\alpha(p)$ for a positive or negative subset coweight
$\alpha$ in $A_0$, since a constant sign can be removed.

We prove that all the orbit sums $\mathsf S_\lambda(p)$ belong to $A_0$
by strong induction on
\[
 \|\lambda\|^2=\sum_{(i,a)\in\Omega}\lambda_{ia}^2.
\]
The coefficient and positive or negative subset-coweight cases have
just been established.
If $\lambda$ has a positive coordinate, set
\[
 f(t)=1_{t>0},\qquad g(t)=t-1_{t>0};
\]
otherwise set
\[
 f(t)=-1_{t<0},\qquad g(t)=t+1_{t<0}.
\]
In either case define $\alpha=f(\lambda)$ and $\gamma=g(\lambda)$
coordinatewise.  Outside the base cases both $\alpha$ and $\gamma$
have squared norm strictly smaller than that of $\lambda$.
Both functions are nondecreasing on $\Z$, and their values have
compatible signs relative to zero.  Thus, for every pair of slots
$d,e$ and every slot $d$,
\[
 (\alpha_d-\alpha_e)(\gamma_d-\gamma_e)\geq0,
 \qquad \alpha_d\gamma_d\geq0.
\]
These statements apply across different colors as well.  Consequently
every exponent in Equation~\eqref{eq:laurent-product} is zero for this
pair, and $r_\alpha r_\gamma=r_\lambda$.

Order the slots at each vertex by their $\lambda$-coordinates.  Then
$\alpha$ and $\gamma$ are ordered in the same way.  The rearrangement
inequality gives, for every relative matching $w\in\mathcal W$,
\[
 \|\alpha+w\gamma\|^2
 \leq\|\alpha+\gamma\|^2=\|\lambda\|^2.
\]
In the equality case, the resulting sum lies in the orbit of $\lambda$.
To see the last assertion directly, sort the matched entries within
each color.  Swapping an inverted pair strictly increases the scalar
product unless one of the two corresponding pairs is tied; swaps
involving ties leave the multiset of sums unchanged.  Hence every
maximal matching has the same multiset of sums as the ordered matching.
All other orbits occurring in a product
$\mathsf S_\alpha(p_\alpha)\mathsf S_\gamma(p_\gamma)$ therefore have
strictly smaller squared norm.

We now verify that the leading coefficient can be an arbitrary
$\mathcal W_\lambda$-invariant dressing.  At the fixed representative
$\lambda$, the decomposition into translates of $\alpha$ and $\gamma$
is unique.  For example, in the positive case, a different assignment
of the $1$-entries of $\alpha$ would assign $1$ to some coordinate
$\lambda_d\leq0$ and $0$ to some coordinate $\lambda_e>0$ of the same
color.  The corresponding $\gamma$-entries would satisfy
$\gamma'_d=\lambda_d-1<\lambda_e=\gamma'_e$, which is a strict
inversion relative to $\alpha'$.  This contradicts maximality.  The
negative case would similarly assign $0$ to a negative coordinate
$\lambda_d$ and $-1$ to a nonnegative coordinate $\lambda_e$ of the
same color; then $\gamma'_d=\lambda_d<\lambda_e+1=\gamma'_e$ is again
a strict inversion relative to $\alpha'$.  Thus the
coefficient of $r_\lambda$ in that product is precisely
$p_\alpha p_\gamma$.

Within each color, the level blocks of $\alpha$ and $\gamma$ are
intervals in the chosen ordering.  Their combined cuts are exactly the
cuts of the level blocks of $\lambda$: two distinct values cannot have
both the same $f$-value and the same $g$-value, since $f(t)+g(t)=t$.
For every $k\geq1$, the power sum on each $\lambda$-block is therefore
a linear combination of power sums on $\alpha$-blocks and
$\gamma$-blocks.  Explicitly, obtain prefix power sums at all cuts
from whichever block partition contains the cut, and subtract
successive prefixes.  Power sums generate each symmetric-polynomial
ring over $\C$, so
\[
 P^{\mathcal W_\lambda}
 =\bigl\langle P^{\mathcal W_\alpha},
                 P^{\mathcal W_\gamma}\bigr\rangle.
\]
As the rings commute, every element of the ring on the right is a
finite sum of products $p_\alpha p_\gamma$ from the two indicated
subrings.  Such products of dressed orbit sums realize any prescribed
leading dressing.  The other terms have polynomial dressings by
Equation~\eqref{eq:laurent-product}; their dressings are invariant under
the stabilizers of their coweights because the full product is
$\mathcal W$-invariant.  The induction hypothesis removes all those
lower terms.  This proves the reverse inclusion, without inverting any
polynomial, and completes the proof.
\end{proof}

\subsection{Multipartitions and cut ideals}

For the rest of this Section, $q$ records half the cohomological degree:
each polynomial variable has cohomological degree two and contributes $q$.
Thus
$\Hilb(M;q)=\sum_d\dim(M^{2d})q^d$ for the evenly graded vector spaces
under consideration.

Let $\mathcal M(\mathbf v)$ be the finite set of arrays
\[
 m=(m_{i,s})_{i\in Q_0,\ s\geq1},\qquad
 m_{i,s}\in\Nbb,\qquad \sum_{s\geq1}s m_{i,s}=v_i.
\]
For $m\in\mathcal M(\mathbf v)$, set
\[
 l_i=\sum_{s\geq1}m_{i,s},\qquad
 \mathbf l=(l_i)_i,\qquad |l|=\sum_i l_i,
\]
and
\[
 D_m=\frac12\sum_{i,j}\sum_{s,t\geq1}
                c_{ij}\min(s,t)m_{i,s}m_{j,t}.
\]
Let $\Omega_m^+$ consist of $m_{i,s}$ distinguished copies of each
type $(i,s)$.  These are the \emph{positive slots}.  For
$d\in\Omega_m^+$ denote its color by $i(d)$ and its size by $s(d)$.
Define
\[
 P_m=\C[x_d:d\in\Omega_m^+],\qquad
 \mathcal W_m=\prod_{i,s}S_{m_{i,s}},\qquad
 b_{de}=c_{i(d)i(e)}\min(s(d),s(e)).
\]
The group $\mathcal W_m$ permutes copies of the same type.  For every
nonempty subset $S\subseteq\Omega_m^+$, put
\begin{equation}\label{eq:cut-polynomial}
 q_S=
 \prod_{d\in S}x_d^{w_{i(d)}}
 \prod_{\substack{d\in S\\ e\in\Omega_m^+\setminus S}}
                    (x_d-x_e)^{b_{de}}.
\end{equation}
Let $J_m\subseteq P_m$ be the ideal generated by these cut polynomials,
and set
\[
 h_m(q)=\Hilb\bigl((P_m/J_m)^{\mathcal W_m};q\bigr).
\]
The ideal is homogeneous and $\mathcal W_m$-stable.  For the empty
array, these conventions give $P_m=\C$, $J_m=0$, and $h_m(q)=1$.

These cut products have a familiar graphical interpretation.
Adjoin a framing vertex $0$, put $w_{i(d)}$ edges between $0$ and each
positive slot $d$, and put $b_{de}$ edges between each unordered pair
of distinct slots $d,e$.
When this multigraph is connected, $J_m$ is the cocircuit ideal of its
reduced incidence configuration, with the row for $0$ deleted.
A minimal disconnecting cut contributes the product of its edge forms,
$x_d$ for framing edges and $x_d-x_e$ for slot edges; every cut product
is divisible by a minimal-cut product.
Such ideals occur in toric hyperk\"ahler cohomology
\cite[Section~7, Equation~(37)]{HauselSturmfels02} and in
permutation-invariant quiver constructions
\cite[Sections~2 and 4.4]{ReinekeRhoadesTewari23}.
A component disjoint from $0$ instead gives a unit cut and zero quotient.
Section~\ref{sec:counting} gives a direct reachability calculation that
retains the $\mathcal W_m$-action and yields the recurrence needed here.

\begin{proposition}\label{prop:filtration-bound}
The quotient $B_0$ for the full topological torus satisfies the
coefficientwise graded bound
\begin{equation}\label{eq:filtration-bound}
 \Hilb(B_0;q)\ \leq
 \sum_{m\in\mathcal M(\mathbf v)}
 q^{C(\mathbf v)-D_m+\mathbf w\cdot(\mathbf v-\mathbf l)}h_m(q).
\end{equation}
Every exponent displayed in the sum is a nonnegative integer.
\end{proposition}

\begin{proof}
Choose, for each color, an identification of its $l_i$ positive slots
with $l_i$ of its $v_i$ gauge slots.  On the remaining $v_i-l_i$ slots
put $s(d)=0$; call these the \emph{zero slots} and write their set as
$\Omega_m^0$.  Define a coweight by
\[
 (\lambda_m)_d=s(d)-1\qquad(d\in\Omega).
\]
The word ``positive'' refers to the part size, so at a rank-three vertex
the sizes $(2,1,0)$ give the coweight $(1,0,-1)$.
The coweight $\lambda_m$ has topological degree zero, since the sum of its coordinates at
vertex $i$ is $\sum_s s m_{i,s}-v_i=0$.  Its stabilizer is
\[
 \mathcal W_{\lambda_m}
 =\mathcal W_m\times\prod_i S_{v_i-l_i}.
\]
The last factors permute the zero slots.

By the degree convention in Proposition~\ref{prop:minuscule}, the
half cohomological degree of $r_{\lambda_m}$ is the polynomial degree
of its Laurent coefficient.  For a remaining arrow $i\to j$, that
degree is
\[
 \sum_{a=1}^{v_i}\sum_{b=1}^{v_j}
       \max(s(ia)-s(jb),0)
 =v_iv_j-\sum_{a=1}^{v_i}\sum_{b=1}^{v_j}
                         \min(s(ia),s(jb)),
\]
because $\sum_a s(ia)=v_i$.  The framing contribution at $i$ is
$w_i\sum_{d\in\Omega_m^+,\ i(d)=i}(s(d)-1)=w_i(v_i-l_i)$.
Summing gives
\begin{equation}\label{eq:multipartition-shift}
 \delta_m:=\deg_q r_{\lambda_m}
 =C(\mathbf v)-D_m+\mathbf w\cdot(\mathbf v-\mathbf l).
\end{equation}
This also proves integrality and nonnegativity.  The formula includes
remaining loops: their ordered pairs of gauge slots give precisely
the diagonal term with $c_{ii}=2n_{ii}$.

For $-1\leq L\leq\sum_i v_i$, define $\mathcal F_{-1}=0$ and, for
$L\geq0$, let $\mathcal F_L\subseteq B_0$ be the sum of the images of
\[
 \mathsf S_{\lambda_m}\bigl(P^{\mathcal W_{\lambda_m}}\bigr),
 \qquad m\in\mathcal M(\mathbf v),\quad |l|\leq L.
\]
Here and below an orbit sum denotes its image in the quotient when it
is viewed in $B_0$.  This is a finite filtration by graded vector
subspaces.  The array $m_{i,1}=v_i$ has $\lambda_m=0$, so the last
stage contains the image of all coefficients.  Those coefficients
generate $B_0$, by Proposition~\ref{prop:minuscule}; hence
$\mathcal F_{\sum_i v_i}=B_0$.  Exhaustion is asserted in the quotient,
and requires no assertion about spanning the entire topological-degree-zero
part of $A_0$ before quotienting.
For later use denote the symbol map of each orbit type by
\[
 \Phi_m:P^{\mathcal W_{\lambda_m}}
 \longrightarrow\mathcal F_{|l|}/\mathcal F_{|l|-1},\qquad
 p\longmapsto[\mathsf S_{\lambda_m}(p)].
\]
It adds the half-degree shift $\delta_m$.

We first eliminate dressings in the zero slots.  For each color $i$
and $k\geq1$, define an ambient Laurent-ring derivation by
\[
 D_{i,k}(p e^\lambda)
 =p e^\lambda\sum_a\lambda_{ia}x_{ia}^k,
 \qquad D_{i,k}(p)=0\quad(p\in\C(x_{ia})).
\]
Additivity of the displayed multiplier in $\lambda$ proves the
Leibniz rule.  This derivation commutes with $\mathcal W$, preserves
the algebra in Equation~\eqref{eq:laurent-model}, and preserves every
topological homogeneous summand.  It therefore preserves the ideal
defining $B_0$.  Its induced derivation on $B_0$ is zero, since it
kills the coefficient generators.  Applying it to an arbitrarily
dressed orbit sum gives
\begin{equation}\label{eq:zero-slot-power-sums}
 \mathsf S_{\lambda_m}\left(
 p\left[
 \sum_{\substack{d\in\Omega_m^+\\i(d)=i}}
                   (s(d)-1)x_d^k
 -\sum_{\substack{z\in\Omega_m^0\\i(z)=i}}x_z^k
 \right]\right)=0\quad\text{in }B_0
\end{equation}
for every $p\in P^{\mathcal W_{\lambda_m}}$.  The invariant dressing
ring is the tensor product of the symmetric-polynomial rings on the
positive type blocks and the zero blocks.  Power sums generate each
zero-block ring over $\C$.  Equation~\eqref{eq:zero-slot-power-sums},
including its arbitrary multiplier $p$, replaces each such power sum
by the indicated positive-slot polynomial.  Thus positive-slot
dressings alone span the image of this orbit type.  These replacements
preserve degree.

Fix now a nonempty $S\subseteq\Omega_m^+$, and split the array of
sizes on all gauge slots as $s=a+b$, by putting $b_d=s(d)$ on $S$ and
$b_d=0$ elsewhere.  Set
\[
 \alpha=a-1,\qquad \beta=b.
\]
The topological degree of $\beta$ is a nonzero vector with
nonnegative entries, so every product
\[
 \mathsf S_\alpha(p_a)\mathsf S_\beta(p_b)
\]
vanishes in $B_0$.  The dressings here may be arbitrary polynomials
invariant under their respective coweight stabilizers, by
Proposition~\ref{prop:laurent-model}.

To analyze the terms of this product, fix $a$ and match it with a
colorwise permutation $b'$ of $b$.  The resulting coweight is $t-1$,
where $t=a+b'$ has nonnegative entries and $\sum_{i(d)=i}t_d=v_i$.
It therefore determines another array in $\mathcal M(\mathbf v)$.
Writing $A=\operatorname{supp}(a)$ and
$B'=\operatorname{supp}(b')$, its length is
\[
 |\operatorname{supp}(t)|
 =|A|+|B'|-|A\cap B'|=|l|-|A\cap B'|.
\]
Every overlap term belongs to $\mathcal F_{|l|-1}$.  For a disjoint
matching, the positive sizes of $t$ are the multiset union of those
of $a$ and $b$ at each color.  Hence its array is exactly $m$; no
other orbit type occurs at length $|l|$.

For the chosen disjoint decomposition, the excess coefficient in
$r_\alpha r_\beta/r_{\lambda_m}$ is $q_S$ up to a nonzero constant
sign.  Here is the calculation for each weight.  For an arrow from
slot $d$ to slot $e$, its excess exponent is
\[
 (a_d-a_e)_+ +(b_d-b_e)_+
                  -(a_d+b_d-a_e-b_e)_+.
\]
This is zero unless one slot belongs to $\operatorname{supp}(a)$ and
the other to $S$.  In that case it is
$\min(s(d),s(e))$.  For the framing weight $-x_d$, the excess is
\[
 (a_d-1)_+ +(b_d)_+ -(a_d+b_d-1)_+,
\]
which equals one on $S$ and zero elsewhere.  These give exactly the
exponents in Equation~\eqref{eq:cut-polynomial}.  With the weight
conventions above, the sign is $(-1)^{\epsilon_S}$, where
\[
 \epsilon_S=
 \sum_{d\in S}w_{i(d)}+
 \sum_{\substack{d\in S\\e\in\Omega_m^+\setminus S}}
 n_{i(d)i(e)}\min(s(d),s(e)).
\]
This integer depends only on the numbers of selected and unselected
copies of each type.  In particular the sign is constant on the
averaging orbit considered next.

It remains to check the full span of the leading dressing relations.
Let
\[
 H=\mathcal W_{\lambda_m},\qquad
 H_S=\mathcal W_\alpha\cap\mathcal W_\beta.
\]
Then $H_S\subseteq H$.  Its blocks are the selected and unselected
copies of each positive type, together with the zero block at every
color.  All disjoint decompositions of the fixed representative
$\lambda_m$ into translates of $\alpha$ and $\beta$ form exactly the
set $H/H_S$.  The coefficient of its leading orbit sum in the product
above is $(-1)^{\epsilon_S}\operatorname{Tr}_{H_S}^H(q_Sp_ap_b)$.
Since the sign is nonzero and all other terms have lower length, we
obtain
\begin{equation}\label{eq:cut-leading-trace}
 \Phi_m\bigl(\operatorname{Tr}_{H_S}^H(q_Sp_ap_b)\bigr)=0,
 \qquad
 \operatorname{Tr}_{H_S}^H(f)=\sum_{h\in H/H_S}h(f).
\end{equation}

Choose $p_a$ to depend only on the positive-size blocks of $a$ and
$p_b$ only on the positive-size blocks of $b$.  Such choices are
allowed by their coweight stabilizers.  Their products span the
$H_S$-invariant polynomial ring on the positive slots, since that ring
is the tensor product of the symmetric-polynomial rings on the
selected and unselected copies of every type.  Moreover $q_S$ is
$H_S$-invariant.  For an arbitrary $f\in P_m$, Reynolds averaging first
over $H_S$ and then over $H$ consequently shows that
\[
 \operatorname{Av}_H(q_S f)
 =\operatorname{Av}_H\bigl(q_S\operatorname{Av}_{H_S}(f)\bigr)
\]
is a linear combination of the relations in
Equation~\eqref{eq:cut-leading-trace}.  All normalizing factors are
nonzero constants in $\C$.  Averaging the generators $q_S$ with
arbitrary multipliers spans $J_m\cap P_m^{\mathcal W_m}$, because
$J_m$ is stable under $\mathcal W_m$ and averaging is exact in
characteristic zero.  Therefore all invariant cut-ideal relations
annihilate the symbol of this orbit type.

Combining the zero-slot elimination with these cut relations gives,
for every $L$, a surjection of graded vector spaces
\begin{equation}\label{eq:filtration-symbol-surjection}
 \bigoplus_{\substack{m\in\mathcal M(\mathbf v)\\|l|=L}}
 (P_m/J_m)^{\mathcal W_m}[\delta_m]
 \twoheadrightarrow \mathcal F_L/\mathcal F_{L-1}.
\end{equation}
Here $[\delta_m]$ means that polynomial half degree $d$ is placed in
half cohomological degree $d+\delta_m$.  Exactness of invariants
identifies the source with the invariant dressing quotient just
obtained.  Summing its Hilbert-series bound over the finite exhaustive
filtration proves Equation~\eqref{eq:filtration-bound}.
\end{proof}

\begin{remark}\label{rem:filtration-edge-cases}
If there are no gauge slots, the unique empty array gives
$A_0=B_0=\C$ and equality in Equation~\eqref{eq:filtration-bound}.
If $\mathbf v\ne0$ and $\mathbf w=0$, the cut on all positive slots
is $q_{\Omega_m^+}=1$, so every summand on the right is zero.
Correspondingly, $e^{\mathbf 1}$ is an invertible monopole of nonzero
topological degree, and $B_0=0$.  The same observation applies to any
unframed connected component.  Full subsets, empty zero blocks, and
rank-one gauge factors require no change in the proof.
\end{remark}

The filtration has reduced the upper bound for $B_0$ to finitely many cut
quotients $(P_m/J_m)^{\mathcal W_m}$.  Section~\ref{sec:counting} evaluates
their generating series and compares the sum in
Equation~\eqref{eq:filtration-bound} with the cohomology at determinant
stability.  No algebra presentation of $B_0$ is inferred from the
vector-space surjections onto its filtration quotients.
